\devoirDS{4}{Chapitres 11 à 13}{Chronophotographie/6,Boule suspendue/7,Lancer vertical/7}
\begin{blocLivre}{ds-4-1}\begin{exoBac}{6}{Chronophotographie}
Les abscisses sont 0 ; 0,20 ; 0,80 ; 1,80 ; 3,20 m, avec 0,20 s entre deux images.
\schema{chronophoto}
\begin{questions}
\item Décrire le mouvement dans le référentiel du laboratoire.
\item Calculer les vitesses approchées en $M_1$ et $M_3$ par les intervalles suivants.
\item Tracer les vecteurs à l’échelle indiquée et conclure sur la compensation des forces.
\end{questions}
\end{exoBac}\end{blocLivre}
\begin{blocLivre}{ds-4-2}\begin{exoBac}{7}{Boule suspendue}
Une boule de 0,200 kg est immobile au bout d’un fil vertical ; $g=10,0$ N/kg.
\schema{suspension}
\begin{questions}
\item Identifier les deux actions principales et les représenter.
\item Calculer le poids et déduire la tension du fil.
\item Ces forces sont-elles une paire d’actions réciproques ?
\item Nommer la force réciproque du poids de la boule.
\end{questions}
\end{exoBac}\end{blocLivre}
\begin{blocLivre}{ds-4-3}\begin{exoBac}{7}{Lancer vertical}
Une balle monte puis redescend ; les actions de l’air sont négligées.
\schema{chute}
\begin{questions}
\item Définir le modèle de chute libre et citer la force retenue.
\item Compléter la figure pour la vitesse et le poids aux trois phases.
\item Pourquoi la vitesse nulle au sommet ne prouve-t-elle pas l’équilibre ?
\item Une balle qui descend à vitesse constante dans un liquide relève-t-elle du même modèle ?
\end{questions}
\end{exoBac}\end{blocLivre}

